% 模板 4 示例：通用研究技术路线图（蓝色主节点+白色子节点+虚线框）
% 编译命令：xelatex demo_roadmap_template4.tex
\documentclass[border=8pt]{standalone}
\usepackage{ctex}
\usepackage{tikz}
\usetikzlibrary{arrows.meta, positioning, calc, shadows}

\begin{document}
\begin{tikzpicture}[scale=1.0,
    main/.style={rectangle, rounded corners=3pt,
        minimum width=5.5cm, minimum height=0.7cm,
        draw=blue!80, line width=0.7pt, fill=blue!6,
        font=\small\bfseries, align=center,
        drop shadow={opacity=0.15, shadow xshift=0.5pt, shadow yshift=-0.5pt}},
    sub/.style={rectangle, rounded corners=2pt,
        minimum width=2.4cm, minimum height=0.6cm,
        draw=teal!70, line width=0.5pt, fill=white,
        font=\footnotesize, align=center},
    dashbox/.style={rectangle, rounded corners=4pt,
        draw=gray!40, dashed, line width=0.7pt, fill=gray!2},
    bigarrow/.style={-stealth, line width=1.4pt, color=blue!70},
    smarrow/.style={-stealth, line width=0.5pt, color=gray!50},
    label/.style={font=\small\bfseries, color=black!70},
]

% ====== 数据预处理 ======
\node[dashbox, minimum width=13cm, minimum height=2cm] (box0) at (0, 0) {};
\node[label, anchor=north west] at (box0.north west) {\scriptsize 数据预处理};
\node[main] (m0) at (0, 0.35) {数据预处理与探索性分析};
\node[sub] (s0a) at (-4.2, -0.65) {异常值检测};
\node[sub] (s0b) at (-1.4, -0.65) {标准化处理};
\node[sub] (s0c) at (1.4, -0.65) {描述性统计};
\node[sub] (s0d) at (4.2, -0.65) {相关性分析};
\foreach \x in {s0a,s0b,s0c,s0d} {\draw[smarrow] (m0) -- (\x);}

\draw[bigarrow] (0, -1.4) -- (0, -2.0);

% ====== 问题一（双层） ======
\node[dashbox, minimum width=13cm, minimum height=4.2cm] (box1) at (0, -4.2) {};
\node[label, anchor=north west] at (box1.north west) {\scriptsize 问题一};
\node[main] (m1) at (0, -2.5) {多维度特征筛选};
\node[sub] (s1a) at (-4.2, -3.4) {Spearman相关};
\node[sub] (s1b) at (-1.4, -3.4) {LASSO回归};
\node[sub] (s1c) at (1.4, -3.4) {随机森林};
\node[sub] (s1d) at (4.2, -3.4) {交集融合$^{\bigstar}$};
\foreach \x in {s1a,s1b,s1c,s1d} {\draw[smarrow] (m1) -- (\x);}

\node[main] (m1b) at (0, -4.7) {贡献度分析};
\node[sub] (s1e) at (-2.8, -5.6) {XGBoost分类$^{\bigstar}$};
\node[sub] (s1f) at (0, -5.6) {SHAP值分解};
\node[sub] (s1g) at (2.8, -5.6) {贡献度排序};
\foreach \x in {s1e,s1f,s1g} {\draw[smarrow] (m1b) -- (\x);}

\draw[bigarrow] (0, -6.3) -- (0, -6.9);

% ====== 问题二（双层） ======
\node[dashbox, minimum width=13cm, minimum height=4.2cm] (box2) at (0, -9.1) {};
\node[label, anchor=north west] at (box2.north west) {\scriptsize 问题二};
\node[main] (m2) at (0, -7.4) {预警模型构建};
\node[sub] (s2a) at (-4.2, -8.3) {XGBoost$^{\bigstar}$};
\node[sub] (s2b) at (-1.4, -8.3) {随机森林};
\node[sub] (s2c) at (1.4, -8.3) {SVM};
\node[sub] (s2d) at (4.2, -8.3) {逻辑回归};
\foreach \x in {s2a,s2b,s2c,s2d} {\draw[smarrow] (m2) -- (\x);}

\node[main] (m2b) at (0, -9.6) {风险分层与规则挖掘};
\node[sub] (s2e) at (-3.6, -10.5) {概率阈值映射};
\node[sub] (s2f) at (-1.0, -10.5) {CART决策树};
\node[sub] (s2g) at (1.8, -10.5) {Apriori关联};
\node[sub] (s2h) at (4.4, -10.5) {核心组合};
\foreach \x in {s2e,s2f,s2g,s2h} {\draw[smarrow] (m2b) -- (\x);}

\draw[bigarrow] (0, -11.2) -- (0, -11.8);

% ====== 问题三（双层） ======
\node[dashbox, minimum width=13cm, minimum height=4.2cm] (box3) at (0, -14.0) {};
\node[label, anchor=north west] at (box3.north west) {\scriptsize 问题三};
\node[main] (m3) at (0, -12.3) {个性化干预方案优化};
\node[sub] (s3a) at (-4.2, -13.2) {MINLP建模};
\node[sub] (s3b) at (-1.4, -13.2) {遗传算法$^{\bigstar}$};
\node[sub] (s3c) at (1.4, -13.2) {枚举验证};
\node[sub] (s3d) at (4.2, -13.2) {敏感性分析};
\foreach \x in {s3a,s3b,s3c,s3d} {\draw[smarrow] (m3) -- (\x);}

\node[main] (m3b) at (0, -14.5) {匹配规律提取};
\node[sub] (s3e) at (-2.8, -15.4) {特征分组};
\node[sub] (s3f) at (0, -15.4) {方案映射};
\node[sub] (s3g) at (2.8, -15.4) {个案输出};
\foreach \x in {s3e,s3f,s3g} {\draw[smarrow] (m3b) -- (\x);}

\end{tikzpicture}
\end{document}
